
% \vspace{-1mm}
\section{Introduction}
%%
Large-scale Vision-Language Pre-training (VLP) models like CLIP \citep{radford2021learning} and ALIGN  \citep{jia2021scaling} have recently demonstrated success across 
various downstream tasks. They learn visual and textual representations from millions of image-text  pairs collected from the Internet and show superior zero-shot ability and robustness. The core technique of these models lies in the global contrastive alignment of the images and texts through a dual-stream model. Such architecture is inference-efficient for downstream tasks like retrieval because the encoders for the two modalities can be decoupled and the image or text representations can be pre-computed offline. 
% But is this global contrastive alignment good enough for all the downstream tasks? Does this kind of multi-modal fusion sufficiently capture the finer-grained interactions? 
% We believe that 
However, CLIP and ALIGN
% one stream of existing works (especially the well-known CLIP and ALIGN) 
model the  cross-modal interaction via solely the similarity of the global feature of each modality, lacking the ability of capturing finer-level information like the relationship between visual objects and textual words.
% (e.g., fine-grained visual objects or textual words). 
In this paper, we develop a simple yet efficient cross-modal finer-grained interaction mechanism for large-scale VLP.

% why late interaction loss
To achieve finer-grained cross-modal interaction, previous methods mainly exploited two kinds of methods.
% VLP frameworks.
(1) One line of work \citep{chen2020uniter,li2020oscar,m5product,li2020unimo,zhang2021vinvl,capture} uses a pre-trained object detector to extract region-of-interest (ROI) features from images, and then fuses it with the paired text through a VLP model.
% Then a VLP model fuses the two modalities by 
% taking the concatenation of the extracted ROI features and the paired text as input. 
This design complicates the pre-training due to  pre-computing and storing  a large number of ROI features. 
In addition, the zero-shot ability of these approaches is usually limited by the predefined number of classes and their performance is also restricted by the quality of the detector.
(2) Another line of work \citep{li2021align, kim2021vilt} 
% \xd{add reference: simVLM, viln etc} 
enforces the token-wise or patch-wise representations from both modalities into the same space and models these finer-grained interactions via cross-attention \citep{li2021align} or self-attention \citep{kim2021vilt}. 
However, these methods are usually less efficient in terms of both training and inference. In particular, during training, cross-attention in \citep{li2021align} requires to be performed in an encoder-decoder structure, while the complexity of the self-attention \citep{kim2021vilt} grows quadratically with the length of the prolonged concatenated sequences of both modalities. During inference, the data from both modalities are intertwined to compute the cross-attention or self-attention, and can not be pre-computed offline as dual-stream models like CLIP and ALIGN.
This can be less efficient for downstream tasks like image/text retrieval and image classification.
% as the encoders can not be decoupled and the image or text representations can not be pre-computed offline like CLIP and ALIGN. 
% \ylw{already mentioned in the above}

In this paper, we propose a large-scale Fine-grained Interactive Language-Image Pre-training framework named FILIP 
to address these limitations.
% , and boost the performance on  more grounded vision and language tasks.
Inspired by \cite{khattab2020colbert}, 
we model  the fine-grained semantic alignment  through a novel cross-modal
late interaction mechanism in the contrastive loss, instead of using cross or self-attention.
Specifically, our fine-grained contrastive learning uses a token-wise maximum similarity between visual  and textual tokens to guide the contrastive objective.
% , while keeping the training and inference efficient.
In this way,
FILIP  successfully leverages the finer-grained expressiveness
among image patches and textual words 
while simultaneously gaining the ability to pre-compute image and text representations offline.
Unlike \cite{khattab2020colbert}, we discard the padded tokens and use average 
instead summation 
of token-wise maximum similarities when computing the image-text alignment,
% between images and texts, 
which enhances the cross-modal representation learning and stabilizes training.
Furthermore, we construct a large-scale pre-training dataset named FILIP300M from the Internet.
% for facilitating the research community. 
Data cleaning and image-text data augmentation are also explored  and proved useful in this work.

% \xd{remember to fill numbers}
Extensive experiments show that by effectively learning fine-grained representations,
% alignment via cross-modal late interaction, 
FILIP achieves state-of-the-art performance on multiple downstream vision-language tasks, including zero-shot image classification and image-text retrieval. For example, FILIP reaches  77.1\% top-1 accuracy for zero-shot ImageNet classification, surpassing  CLIP with less training data.
Visualizations on  word-patch alignment further show that FILIP learns meaningful finer-grained features with promising localization ability. 

% On XX and XX, FILIP achieves absolute improvements of XX\% and XX\% compared to XX method. On zero-shot classification and image-text retrieval, our FILIP outperforms current state-of-the-art visual-language pre-training methods.

%To support more downstream generation tasks such as image captioning, FILIP further incorporates a light interaction layer on top of the original dual-encoder architecture. 



%In the domain of natural language processing, a late interaction loss is proposed in \citep{khattab2020colbert} to model the fine-grained similarity between the   independently encoded query and the document representations from  the BERT model, in order to solve the information retrieval task.
%Inspired by \cite{khattab2020colbert}, we can delay the interaction without cross-attention or self-attention and yet retain the fine-grained interaction between text and image, through a similar late interaction loss. \footnote{\hou{add description of the functionality of late interaction loss}}
%In this way, we can leverage the finer-grained expressiveness while simultaneously gaining the ability to pre-compute image and text representations offline. 
%Unlike the original late interaction loss in \cite{khattab2020colbert}, we do not consider the padded textual tokens when computing the token-wise similarities to avoid letting them mislead the model to align image patches to these padded tokens. In addition, we use average the token-wise maximum similarities instead of original summarization to stabilize training. 